
\paragraph{Beyond exponential memory.} The proof of part~(iii) uses only one property of the feedback traders: their aggregate flow is a fixed linear function of past public signals. So the result does not depend on exponential memory.

\begin{proposition}[Any linear population]\label{prop:anylinear}
In the market of Theorem~\ref{thm:eqinv}, replace the cohorts by any finite population of feedback traders $j=1,\ldots,J$ whose flows are causal linear functions of the public signal,
\[
 f_{j,t}=\sum_{m<t}\kappa^{(j)}_m\,s_{t-1-m},\qquad \kappa^{(j)}=(\kappa^{(j)}_0,\kappa^{(j)}_1,\ldots)\ \text{arbitrary real sequences},
\]
and write $k=\sum_j\kappa^{(j)}$ for the population's kernel.
\begin{enumerate}[label=(\roman*),itemsep=1pt]
\item Two populations produce the same aggregate flow on dates $t\le T$ for every signal sequence if and only if their kernels agree for $m<T$.
\item If their kernels agree for $m<T$, then for every strategy profile and every realization the price and order-flow paths, the insider's trades and realized profit, and the dealer's inventory coincide on $t\le T$, and the two games have the same equilibria under any equilibrium concept defined from payoffs and the public history.
\item Any decomposition of a kernel into individual kernels is consistent with the same outcomes. Prices, flows and all equilibrium outcomes therefore carry no information about the decomposition beyond $k$.
\end{enumerate}
\end{proposition}

\begin{proof}
(i) Aggregate flow is $f_t=\sum_{m<t}k_ms_{t-1-m}$. If the kernels agree for $m<T$, the flows agree on $t\le T$ for every signal. Conversely, the impulse signal $s_j=\mathbf 1\{j=0\}$ gives $f_{m+1}=k_m$. (ii) The strong induction in the proof of Theorem~\ref{thm:eqinv}(iii) (Supplement~\ref{app:eqinv}) uses only that $f_t$ is the same function of the signals before $t$ in both markets, which is part~(i). (iii) follows from (ii), since outcomes depend on the population only through $k$.
\end{proof}

\paragraph{What the exponential case adds.} For general linear traders, nothing about composition beyond $k$ can be learned from anything that prices determine, because any split of $k$ is observationally equivalent. Exponential memory is the case in which more can be learned, and the rest of the paper prices it. Then $k_m=\int r^m\,dM(r)$ is a moment sequence, and the kernel determines the merged measure $M$ (Theorem~\ref{thm:eqinv}(ii) on an unbounded horizon). A finite record determines only finitely many moments, which is why separating close groups is costly. Proposition~\ref{prop:anylinear} is machine-checked for arbitrary kernels (\texttt{p1\_anylinear\_iff}, \texttt{p1\_anylinear\_invariance}); the passage to equilibria is on paper, as in Theorem~\ref{thm:eqinv}.

\paragraph{The analogy with aggregate demand.} The kernel is the response of aggregate feedback demand to past prices. \citet{auclert2024} show that the analogous object for households, the matrix of intertemporal marginal propensities to consume, is a sufficient statistic for the transmission of fiscal policy. In the same way, $k$ summarizes every heterogeneity among feedback traders that matters for prices. The merged measure is the distribution of individual autoregressive coefficients that governs an aggregate in \citet{granger1980}. Two points distinguish the market version. Sufficiency holds inside a game with an insider and a dealer, not only for the law of an aggregate series. And as part~(iv) of Theorem~\ref{thm:eqinv} shows, it fails for group-level quantities, which matter for targeted supervision (Section~\ref{sec:supervision}).

\subsection{What prices identify}

Theorem~\ref{thm:eqinv} says that populations with the same kernel cannot be told apart by anything that price-based trading determines. The converse is that prices determine the kernel. In the probing experiment of Section~\ref{sec:model}, the response of prices to trades is a fixed function of $(a,\kappa,k)$, and it determines the merged measure.

